%============================================================ 1. 서론
\section{서론}
\label{sec:intro}

\subsection{배경}
\begin{figure*}[tp]
\centering
\includegraphics[width=0.9\textwidth]{fig1.png}
\caption{제안 체계의 개요 --- 전략 계층(위)의 무리별 배정과 기동 계층(아래)의 그물벽 위치 결정}
\label{fig:overview}
\end{figure*}

다수의 저비용 무인수상정(Unmanned Surface Vehicle, USV)이 고가치 함정 한 척을 향해 여러
방향에서 동시에 돌진하는 자폭 공격은 대표적인 비대칭 위협이다. 유한한 방어 자산을 다수
표적에 배분하는 문제는 무기-표적 할당(weapon-target assignment)으로 오래 연구되어
왔으나~\citep{ahuja2007wta}, 표적 수가 요격 자산을 넘어서면 최적 할당으로도 처리되지 않는
표적이 남는다.

무인 위협에 대한 기존 대응은 두 갈래로 구분된다~\citep{wang2021cuas}. 표적의 통신·항법
·센서를 교란하는 연성 대응(soft-kill)은 위협 자체를 제거하지 못하며, 전파를 방사하지 않는
유선 조종 표적에는 작용하지 않는다~\citep{sutea2026fiber}. 표적을 직접 타격하는 경성
대응(hard-kill)은 위협을 확실히 제거하지만, 수천 달러 규모의 표적에 발당 수십만에서
수백만 달러의 요격 유도탄을 소모하는 비용 구조를 지탱하기 어렵고~\citep{rumbaugh2024cost},
한 번에 한 표적만 처리하므로 표적 수가 교전 용량을 넘어서면 포화가 발생한다.

본 연구는 대안으로 저비용 그물을 이용한 물리적 포획을 채택한다. 이는 무인기 대응 분야에서
독립된 대응 범주로 다루어져 왔으며~\citep{wang2021cuas}, 표적의 통신 방식과 무관하게
작용하고 표적을 정지시킨다는 점에서 앞의 두 한계를 비켜간다. 방어정은 적선보다 느리므로
추격으로 따라잡을 수 없고, 대신 적선이 지나갈 접근 회랑을 앞질러 그물벽으로 차단한다
(Fig.~\ref{fig:overview}). 그림의 함정 형상과 패널은 개념 전달용 삽화이며, 실제 입출력 형식은 \S\ref{sec:method}에 정의한다.

\subsection{문제의식}

이 방어 문제는 성격이 다른 두 판단을 동시에 요구한다. 하나는 \emph{어느 방어정이 어느
적 무리를 맡을 것인가}라는 이산적·전술적 배정이고, 다른 하나는 \emph{맡은 무리를 막기 위해
해면의 어느 지점에 그물벽을 놓을 것인가}라는 연속적·기하적 기동이다. 전자는 전장 전체의
맥락과 우선순위 판단을 요구하고, 후자는 국소적이지만 매 결정마다 새로 계산되어야 한다.
두 판단은 같은 주기로 갱신되지만 \emph{산출에 걸리는 시간이 크게 다르다}. 본 연구는 느린 쪽인 배정을 언어모델 지휘관에게, 빠른 쪽인 기동을 학습된 정책에 맡긴다(Fig.~\ref{fig:overview}). 운용 제약도 설계에 들어간다 --- 결정 주기는 25초, 온보드 자원은 단일 GPU이며, 방어정은 적선보다 느리고 그물은 한 척당 몇 장뿐이다. 따라서 묻는 것은 ``언어모델과 학습 정책을 쓸 수 있는가''가 아니라 \emph{그 제약 아래에서 각 계층이 실제로 무엇을 얼마나 사는가}이다.


\subsection{배치를 전제로 한 설계}

본 연구는 시뮬레이션 안에서만 성립하는 체계를 목표로 하지 않았다. 방법을 고를 때마다
\emph{온보드에서 돌 수 있는가}를 함께 물었고, 그 선택들이 본문 곳곳에 흩어져 있으므로
여기에 모아 둔다.

\begin{itemize}[leftmargin=*, itemsep=2pt, topsep=3pt]
\item \textbf{실해역 지형에서 학습·평가한다.} 교전 해역은 가상의 빈 바다가 아니라
      OpenStreetMap 의 실제 해안선(통영 매물도 인근)을 래스터화한 것이며, 육지는 충돌
      판정과 관측 채널 양쪽에 들어간다(\S\ref{sec:obs}).
\item \textbf{전장을 이미지로 보고, 행동을 픽셀로 낸다.} 관측은 선박별 벡터의 나열이
      아니라 모선을 중앙에 둔 $50\times50$ 격자 위의 15채널 영상이고, 정책은 그 위에서
      픽셀을 지목하는 U-Net 구조이다. 행동은 연속 좌표의 회귀가 아니라 격자의 유효 픽셀
      가운데 하나를 고르는 \emph{이산적인 선택}이므로, 정책은 해면 전체에 대한 점수맵을
      낸다. 연속 좌표 회귀에서는 불가능 영역을 피하도록 학습하거나 사후에 투영해야 하지만,
      이산 선택에서는 그런 자리가 애초에 뽑히지 않게 막을 수 있다. 선박 수가 바뀌어도
      입력·출력의 형태가 변하지 않으며, 배치 통계에 의존하는 정규화를 쓰지 않아 배치 크기
      1의 온보드 추론이 학습과 같은 연산을 수행한다(\S\ref{sec:obs}, \S\ref{sec:deploy}).
\item \textbf{언어모델은 느린 판단만 맡고, 빠른 판단을 막지 않는다.} 배정은 언어모델
      지휘관이, 기동은 학습된 정책이 맡는다. 둘은 비동기로 분리되어 지휘관의 응답을
      기다리는 동안에도 기동은 멈추지 않는다 --- 응답 속도가 크게 다른 두 판단 계층을
      같은 결정 주기 안에서 결합한 것이다.
\item \textbf{출력이 곧 항법 명령이다.} 정책의 출력은 제어 신호가 아니라 경유점 두 개이며,
      이 형식은 특정 선체 동역학에 묶이지 않는다. 저수준 추종은 학습과 분리되어 있어
      실제 선체의 동역학으로 교체할 수 있고, 그만큼 시뮬레이션에서 실기로 옮기기 쉽다
      (\S\ref{sec:deploy}).
\end{itemize}

\subsection{기여}

본 논문의 기여는 \emph{무엇이 새로 가능해졌는가}로 정리된다.

\begin{enumerate}[leftmargin=*, itemsep=3pt, topsep=4pt]
\item \textbf{공학 해법.} 적보다 느린 방어정 세 척으로 자폭 군집을 막는 2계층 방어 체계 ---
      배정은 언어모델 지휘관, 그물벽 위치는 15채널 격자 관측 위의 U-Net 점수맵 정책이 맡고,
      느린 계층만 비동기로 분리한다. 관측·행동이 격자 위에 있어 방어정·적선 수가 바뀌어도
      같은 모델(244{,}389 파라미터)이 돌고, 배치 통계 없는 정규화와 경유점 출력으로 온보드
      실시간 추론과 항법 스택 연결이 그대로 된다(\S\ref{sec:deploy}).
\item \textbf{검증 결과.} 적보다 느린 방어정 세 척이 그물벽만으로 자폭 군집을 같은 수준으로 막으면서 포획당 그물 소모를 44.5\,\%, 총 선회량을 46.6\,\% 줄이고, 언어모델 지휘관은 능력 임계 위의 모델에서만 이득이며 로컬 7B급은 오히려 손해라는 것을 --- 조합 최적화 후처리 없이 시드 대응 $2\times2$ 설계로 두 계층의 기여를 분리 측정하여 720 에피소드에서 보였다. 기동 계층 단독으로도 포획당 그물 22.9\,\%, 그물 걸림
      24.3\,\% 를 줄이며, 이 조건은 외부 API 없이 재현된다. 10개 지표 중 9개에서 두 계층의
      기여는 가산적이다 --- 서로를 돕지도 잠식하지도 않는다.
\item \textbf{배치 판단 기준.} 언어모델 지휘관은 배정 커버리지가 아니라 \emph{재계획 간 계획
      안정성}으로 평가해야 하며(가장 많이 배정한 모델이 가장 나빴다), 로컬 중형 모델의 응답
      지연은 결정 주기를 넘길 수 있어(14B 호출의 4.2\,\%) 로컬 배치에서는 비동기 분리가 전제다.
      언제 지휘관 계층을 켜고 어느 모델을 쓸지에 대한 판단 기준을 수치로 제공한다.
\end{enumerate}

\noindent
논문의 구성은 다음과 같다. \S\ref{sec:related}에서 관련 연구를 정리하고,
\S\ref{sec:problem}에서 문제와 평가 지표를 정의한다. \S\ref{sec:method}에서 제안 체계를
기술하고, \S\ref{sec:experiment}에서 실험 설계와 통계 절차를 밝힌다.
\S\ref{sec:results}에서 결과를 보고하고, \S\ref{sec:discussion}에서 해석과 한계,
결론을 제시한다.
